\scanpage{42}
% Pencil archive foliation: 41.
\uline{11} \uline{Interprétation topologique du type}
\deleted{\unsure{filtre}} \uline{dimensionnel} (Cas ``géométrique'').

Soit $X$ \deleted{projectif} non singulier sur $k$
alg. clos. Considérons la filtration $X$
par la codimension, et la suite spectrale
\[
H^n(X,\unsure{\mathbb{Q}|\mathbb{Z}}_\ell)
\Longleftarrow
\inserted{E_1^{pq}=}
\coprod_{x\in X[p]}H^{q-p}(x,\unsure{\mathbb{Q}|\mathbb{Z}}_\ell(-p))
\deleted{\unsure{\xi}}
\]
\margin{Attention aux \unsure{infinis} $E_1$ et $E_2$ !!}
où on pose
\[
\left\{\begin{array}{l}
X[p]=\{x\in X\mid\dim\mathcal{O}_{X,x}=p\}\\
H^\alpha(x,\mathbb{Z}_\ell(-p))=
\underset{\substack{\longrightarrow\\
U\text{ ouvert }\ne\varnothing\text{ de }\bar{x}}}{\lim}
H^\alpha(U,\unsure{\mathbb{Q}|\mathbb{Z}}_\ell(\unsure{-p|p})))
\end{array}\right.
\]
\figure{\begin{tikzpicture}
\draw (-.4,0)--(3.2,0);
\draw (0,-.4)--(0,2.6);
\draw (0,0)--(3.2,2.5);
\node at (3.6,0) {$p$};
\node at (-.2,2.7) {$q$};
\node at (2.7,1.1) {$0$};
\end{tikzpicture}}
On veut que \deleted{la \unsure{provenance de}} cette suite spectrale
provienne d'une suite spectrale de
motifs (du moins à partir de $E_r^{pq}$ avec
$r\geq2$ ; sinon il faudrait parler de
ind-motifs, ou bien prendre une
filtration \uline{finie} convenable de $X$ par des
sous-schémas fermés].
\margin{\unsure{codim $p$|en dim $p$}}
Le morceau
de la filtration $\geq p$ est visiblement de
type dimensionnel $\leq n-2p$. On veut que
ce soit \uline{exactement} le morceau de type
dimensionnel $n-p$.
\margin{si $X$ est \uuline{\uline{projectif}}
(\unsure{l'hypothèse essentielle}, \unsure{penser}
à $X$ de dim $1$ !)}

\figure{\begin{tikzpicture}
\draw (0,1.2)--(0,.15) .. controls (-.16,.05) and (-.16,-.13) .. (0,-.08) .. controls (.12,-.03) and (.13,.13) .. (0,.15);
\draw (6,1.2)--(5.95,.15) .. controls (5.78,.07) and (5.78,-.10) .. (5.94,-.07) .. controls (6.09,-.02) and (6.09,.13) .. (5.95,.15);
\end{tikzpicture}}
